% Part: first-order-logic
% Chapter: proof-systems
% Section: axiomatic-deduction

\documentclass[../../../include/open-logic-section]{subfiles}

\begin{document}

\iftag{FOL}
      {\olfileid[tr]{fol}{prf}{axd}}
      {\olfileid[tr]{pl}{prf}{axd}}

\olsection{Aksiyomatik \usetoken{P}{derivation}}

Aksiyomatik !!{derivation}s kullanan en eski ve en basit mantıksal sistemler,
aksiyomatik !!{derivation} sistemleridir. Bu sistemlerde !!{derivation}s
yalnızca !!{sentence}s dizileridir. Bir !!{sentence}s dizisi, doğru bir
!!{derivation} sayılmak için dizideki her !!{sentence}~$!A$ bakımından
aşağıdaki koşullardan birini sağlamalıdır:
\begin{enumerate}
\item $!A$ bir aksiyomdur veya
\item $!A$, verilmiş bir kümede !!a{element} olarak yer alır; bu küme $\Gamma$ ile
  gösterilen !!{sentence}s kümesidir ya da
\item $!A$ bir çıkarım kuralıyla gerekçelendirilmiştir.
\end{enumerate}
$!A$'nın aksiyom olması için, sabit sayıdaki !!{sentence} şemalarından birinin
biçiminde olması gerekir. \iftag{FOL}{Birinci dereceden mantık}{Önerme mantığı}
için yeterli (sağlam ve
tam) bir !!{derivation} sistemi sağlayan birçok aksiyom şeması kümesi vardır.
Bunların bazıları yönettikleri bağlaçlara göre düzenlenmiştir; örneğin
\[
!A \lif (!B \lif !A) \qquad !B \lif (!B \lor !C) \qquad (!B \land !C) \lif !B
\]
şemaları $\lif$, $\lor$ ve~$\land$ bağlaçlarını yöneten yaygın aksiyomlardır.
Bazı aksiyom sistemleri asgari sayıda aksiyom kullanmayı amaçlar. İlkel kabul
edilen bağlaçlara bağlı olarak, tek bir aksiyomdan oluşan aksiyom sistemleri
bulmak bile mümkündür.

Çıkarım kuralı, !!a{sentence} için, !!a{derivation} içinde
gerekçelendirilmesini sağlayan yeterli koşulu veren koşullu bir ifadedir.
Modus ponens bu tür çok yaygın
bir kuraldır: $!A$ ile $!A \lif !B$ zaten gerekçelendirilmişse $!B$'nin de
gerekçelendirilmiş olduğunu söyler. Bu, !!a{derivation} içinde
!!{sentence}~$!B$'yi içeren satırın, hem $!A$ hem de $!A \lif !B$ (bir
!!{sentence}~$!A$ için) !!{derivation} içinde $!B$'den önce yer alıyorsa
gerekçelendirilmiş olduğu anlamına gelir.

Aksiyomatik !!{derivation}s esas alınarak $\Proves$ bağıntısı şöyle
tanımlanır: !!a{derivation} varsa, son satırında !!{sentence}~$!A$ yer alıyorsa
ve (2) ile gerekçelendirilen !!{sentence}s bu !!{derivation} içinde
$\Gamma$'dan seçilmişse $\Gamma \Proves !A$ olur. Eşdeğer olarak, kullanılan
varsayımların kümesi $\Gamma$'nın sonlu bir alt kümesidir. $\Gamma$ boşken
$!A$ için !!a{derivation} varsa, yani her !!{sentence} bu !!{derivation}
içinde (1) veya (3) ile gerekçelendiriliyorsa $!A$ bir teoremdir. Örneğin
aşağıdaki !!a{derivation}, $\Proves !A \lif (!B \lif (!B \lor !A))$ olduğunu gösterir:
\begin{derivation}
  1. & $!B \lif (!B \lor !A)$ \\
  2. & $(!B \lif (!B \lor !A)) \lif (!A \lif (!B \lif (!B \lor !A)))$\\
  3. & $!A \lif (!B \lif (!B \lor !A))$
\end{derivation}
1. satırdaki !!{sentence}, $!A \lif (!A \lor !B)$ aksiyomu biçimindedir
($!A$ ile $!B$'nin rolleri ters çevrilmiştir). 2. satırdaki tümce ise $!A
\lif (!B \lif !A)$ aksiyomu biçimindedir. Dolayısıyla iki satır da
gerekçelendirilmiştir. 3. satır modus ponens ile gerekçelendirilir: bu satırı
$!D$ ile kısaltırsak 2. satır $!C \lif !D$ biçimindedir; burada $!C$, $!B
\lif (!B \lor !A)$, yani 1. satırdır.

$\Gamma \Proves \lfalse$ ise $\Gamma$ kümesi tutarsızdır. Tam bir aksiyom
sistemi her $!A$ için $\lfalse \lif !A$ ifadesini de ispatlar; dolayısıyla
$\Gamma$ tutarsızsa her $!A$ için $\Gamma \Proves !A$ olur.

Mantık için aksiyomatik !!{derivation}s veren ilk kişi Gottlob Frege'ydi; bu
nedenle çoğu zaman modern mantığın ilk eseri sayılan 1879 tarihli
\emph{Begriffsschrift}'inde verdi. Bu sistemler Alfred North Whitehead ile
Bertrand Russell'ın \emph{Principia Mathematica}'sında ve 1920'lerde David
Hilbert ile öğrencilerinin çalışmalarında yetkinleştirildi. Bu yüzden çoğu
zaman “Frege sistemleri” veya “Hilbert sistemleri” diye adlandırılırlar. Bir
mantık için aksiyomatik sistem bulmak çoğu zaman kolay olduğundan son derece
çok yönlüdürler. !!{derivation}s çok basit bir yapıya ve yalnızca bir veya iki
çıkarım kuralına sahip olduğundan, onlar \emph{hakkında} sonuçlar ispatlamak
da görece kolaydır. Ancak uygulamada kullanımları çok zordur; yani ispatları
bulmak ve yazmak güçtür.

\end{document}
